\documentclass[11pt,a4paper]{article}

% ====== Packages ======
\usepackage[utf8]{inputenc}
\usepackage[T1]{fontenc}
\usepackage[margin=1in]{geometry}
\usepackage{graphicx}
\usepackage{amsmath,amssymb}
\usepackage{booktabs}
\usepackage{longtable}
\usepackage{array}
\usepackage{hyperref}
\usepackage{xcolor}
\usepackage{listings}
\usepackage{enumitem}
\usepackage{titlesec}
\usepackage{tikz}
\usetikzlibrary{shapes.geometric, arrows.meta, positioning, fit, backgrounds}
\usepackage{caption}
\usepackage{subcaption}

% ====== Chinese support (CJK via xeCJK if compiled with XeLaTeX) ======
% If compiled with XeLaTeX or LuaLaTeX, uncomment the following:
% \usepackage{xeCJK}
% \setCJKmainfont{Noto Serif CJK SC}

% ====== Hyperref styling ======
\hypersetup{
  colorlinks=true,
  linkcolor=black,
  urlcolor=blue!60!black,
  citecolor=blue!60!black,
  pdftitle={Soma: A Knowledge-Augmented, Personalized Therapeutic Music System},
  pdfauthor={Soma Team}
}

% ====== Code listing style ======
\definecolor{codebg}{RGB}{248,248,248}
\definecolor{codekw}{RGB}{0,0,180}
\definecolor{codestr}{RGB}{160,0,0}
\definecolor{codecmt}{RGB}{0,120,0}
\lstset{
  backgroundcolor=\color{codebg},
  basicstyle=\ttfamily\footnotesize,
  keywordstyle=\color{codekw}\bfseries,
  stringstyle=\color{codestr},
  commentstyle=\color{codecmt}\itshape,
  showstringspaces=false,
  breaklines=true,
  frame=single,
  rulecolor=\color{black!20},
  numbers=left,
  numberstyle=\tiny\color{gray},
  numbersep=6pt,
  captionpos=b
}

% ====== Title ======
\title{
  \vspace{-1.5cm}
  \textbf{Soma: A Knowledge-Augmented, Personalized\\ Therapeutic Music Generation System}\\[0.4em]
  \large Mid-Term to Final Engineering Report
}
\author{Soma Project Team}
\date{May 2026}

% ====== Section format ======
\titleformat{\section}{\Large\bfseries}{\thesection}{1em}{}
\titleformat{\subsection}{\large\bfseries}{\thesubsection}{1em}{}

\begin{document}
\maketitle

\begin{abstract}
\noindent
Soma is a wellness-oriented system that turns biometric signals, user context,
and music preferences into personalized, AI-generated therapeutic music
sessions. This report documents the engineering progress made between the
mid-term checkpoint (commit \texttt{32a3159}) and the final submission, covering
51 commits, 59 modified files, and approximately 17{,}100 added / 8{,}100
deleted lines of code. The work spans (i) a full project re-organization into a
standard multi-package layout, (ii) a substantially rewritten music-AI engine
with a continuous arousal model and per-user personalization, (iii) a from-scratch
clinical music knowledge base built on a GraphRAG architecture with hybrid
dense--lexical retrieval, (iv) a complete reconstruction of the web client from
a single HTML file into a multi-workspace single-page application, and
(v) groundwork for an iOS / HealthKit migration. We describe the system
architecture, the algorithmic and engineering decisions behind each subsystem,
the test coverage we introduced, and the limitations and future directions of
the project.
\end{abstract}

\tableofcontents
\vspace{1em}
\hrule
\vspace{1em}

% ============================================================
\section{Introduction}
% ============================================================

\subsection{Background and Motivation}

Music has long been used as a non-pharmacological aid for anxiety regulation,
sleep onset, focus, grounding, and energy modulation. Recent advances in
generative audio models (e.g.\ Suno, MusicGen) make it feasible to synthesize
arbitrary therapeutic music on demand, but raw generative capability alone is
not sufficient for a wellness product. Three gaps remain:

\begin{enumerate}[leftmargin=*]
  \item \textbf{Personalization gap.} Off-the-shelf generation does not take
    into account a user's physiological state, history, or preference profile.
  \item \textbf{Evidence gap.} Recommendations should be traceable to clinical
    or musicological evidence, not produced by an opaque prompt template.
  \item \textbf{Experience gap.} A user arriving in an anxious or fatigued
    state has limited patience for technical setup; the interface must reduce
    cognitive load, not add to it.
\end{enumerate}

\subsection{Goal of This Reporting Period}

The mid-term version of Soma was a working prototype: a single-file web page
plus a notebook-driven Python pipeline. The objective for the final phase was
to evolve it into \textbf{a deliverable system}: structured, modular, tested,
and grounded in retrievable evidence. Concretely, we set five goals:

\begin{enumerate}[leftmargin=*]
  \item Re-organize the repository into a maintainable multi-package layout.
  \item Upgrade the music-AI engine from discrete rule mappings to a
    continuous, personalized strategy.
  \item Build a clinical music knowledge base with GraphRAG-style retrieval.
  \item Reconstruct the web client into a coherent multi-workspace SPA.
  \item Establish a test suite and write the foundation for iOS migration.
\end{enumerate}

\subsection{Quantitative Summary}

\begin{table}[h]
\centering
\caption{Aggregate change statistics, mid-term to final.}
\begin{tabular}{lr}
\toprule
\textbf{Metric} & \textbf{Value}\\
\midrule
Commits                & 51 \\
Files changed          & 59 \\
Lines added            & $\approx 17{,}100$ \\
Lines deleted          & $\approx 8{,}100$ \\
New Python sub-packages & 2 (\texttt{apps/api}, \texttt{music\_ai\_module/knowledge}) \\
New test files         & 12 ($\approx 800$ LOC) \\
New documentation files & 6 (incl.\ a 1{,}782-line iOS migration plan) \\
\bottomrule
\end{tabular}
\end{table}

% ============================================================
\section{System Architecture}
% ============================================================

Figure~\ref{fig:arch} shows the overall architecture after the final-phase
refactor. The system is organized into four horizontal layers: \emph{Client},
\emph{API}, \emph{Music-AI Engine}, and \emph{Knowledge \& Data}.

\begin{figure}[h]
\centering
\begin{tikzpicture}[
  node distance=0.6cm and 0.6cm,
  layer/.style={draw, rounded corners, fill=blue!5, minimum width=14cm, minimum height=1.6cm, align=center},
  box/.style={draw, rounded corners, fill=white, minimum width=3cm, minimum height=0.9cm, align=center, font=\small},
  arrow/.style={-{Stealth[length=2mm]}, thick}
]
  % Client layer
  \node[layer] (client) {};
  \node[above=0.05cm of client.north west, anchor=south west] {\textbf{Client (Web SPA, future iOS)}};
  \node[box] at ($(client.west)+(2.2,0)$) (sidebar) {Sidebar Shell};
  \node[box] at ($(client.west)+(5.8,0)$) (sess)    {Session\\(Check-in/Playback)};
  \node[box] at ($(client.west)+(9.4,0)$) (rep)     {Diagnostic\\Report};
  \node[box] at ($(client.west)+(12.6,0)$) (set)    {Settings /\\History};

  % API layer
  \node[layer, below=of client] (api) {};
  \node[above=0.05cm of api.north west, anchor=south west] {\textbf{API Layer (FastAPI, \texttt{apps/api/main.py})}};
  \node[box] at ($(api.west)+(2.2,0)$) (sapi)  {/session};
  \node[box] at ($(api.west)+(5.8,0)$) (fapi)  {/feedback,\\/history};
  \node[box] at ($(api.west)+(9.4,0)$) (hkapi) {/healthkit/\\mock/run};
  \node[box] at ($(api.west)+(12.6,0)$) (kapi) {/knowledge\\(retrieval)};

  % Engine layer
  \node[layer, below=of api, fill=orange!5] (eng) {};
  \node[above=0.05cm of eng.north west, anchor=south west] {\textbf{Music-AI Engine (\texttt{music\_ai\_module/})}};
  \node[box] at ($(eng.west)+(2.2,0)$) (proc) {Processor\\(physio.\ pipeline)};
  \node[box] at ($(eng.west)+(5.8,0)$) (pers) {Personalization\\\& Strategy};
  \node[box] at ($(eng.west)+(9.4,0)$) (comp) {Compiler\\(prompt builder)};
  \node[box] at ($(eng.west)+(12.6,0)$) (pipe) {Pipeline\\(orchestration)};

  % Knowledge layer
  \node[layer, below=of eng, fill=green!5] (kn) {};
  \node[above=0.05cm of kn.north west, anchor=south west] {\textbf{Knowledge \& Data (GraphRAG)}};
  \node[box] at ($(kn.west)+(2.2,0)$) (extr)  {Extractor /\\Fetcher};
  \node[box] at ($(kn.west)+(5.8,0)$) (graph) {Graph Store\\+ Embeddings};
  \node[box] at ($(kn.west)+(9.4,0)$) (retr)  {Hybrid\\Retriever};
  \node[box] at ($(kn.west)+(12.6,0)$) (audit){Clinical\\Auditor};

  % Arrows
  \draw[arrow] (client.south) -- (api.north);
  \draw[arrow] (api.south)    -- (eng.north);
  \draw[arrow] (eng.south)    -- (kn.north);
  \draw[arrow, dashed] (kn.east) to[bend right=25] node[right,font=\scriptsize]{evidence} (eng.east);
\end{tikzpicture}
\caption{Soma final-phase architecture. The client communicates with a FastAPI
service that orchestrates the music-AI engine, which in turn queries the
GraphRAG knowledge base for evidence-grounded musical parameters.}
\label{fig:arch}
\end{figure}

\subsection{Repository Layout}

\begin{lstlisting}[language=bash, caption={Final-phase repository layout.}]
Soma/
├── apps/
│   ├── api/main.py            # FastAPI service (~304 LOC)
│   └── web/index.html         # Sidebar SPA (~9,866 LOC)
├── music_ai_module/
│   ├── processor.py           # Physiological pipeline
│   ├── personalization.py     # Per-user strategy
│   ├── compiler.py            # Prompt / parameter compiler
│   ├── pipeline.py            # End-to-end orchestration
│   ├── style_maps.py          # Genre / style mappings
│   └── knowledge/             # NEW: GraphRAG sub-package
│       ├── schemas.py, extractor.py, fetcher.py,
│       ├── embeddings.py, graph_store.py,
│       ├── retriever.py, auditor.py, cli.py
├── data/
│   ├── case.md, case_audio_urls.json
│   └── knowledge/             # Seed graph + chunks + sources
├── docs/                      # PRD, API contract, demo script,
│                              # iOS migration plan (1,782 LOC)
├── tests/                     # 12 test modules, ~800 LOC
├── scripts/                   # api_test.js, generate_cases.js
├── notebooks/pipeline_demo.ipynb
├── pyproject.toml, requirements.txt, .env.example
└── README.md, PRODUCT.md
\end{lstlisting}

% ============================================================
\section{Engineering Foundation: Repository Refactor}
% ============================================================

The mid-term codebase had several structural problems: a 3{,}834-line
\texttt{web.html} and a 3{,}675-line \texttt{Music AI.ipynb} sat at the
repository root, scripts and JSON fixtures were unsorted, and there was no
unified dependency or environment manifest. The first phase of the final-period
work was to address these problems before any feature work, because the
downstream goals (modular engine, knowledge subsystem, test coverage) were not
viable on the old layout.

The refactor performed four actions:

\begin{enumerate}[leftmargin=*]
  \item \textbf{Standard Python layout.} Created \texttt{apps/}, \texttt{scripts/},
    \texttt{data/}, \texttt{docs/}, \texttt{notebooks/}, and \texttt{tests/}.
    Each Python directory received an \texttt{\_\_init\_\_.py} so that the
    project is importable as a package and the test suite can use absolute
    imports.
  \item \textbf{Removed root-level monoliths.} The notebook and the single-file
    web page were deleted from the root and replaced by their refactored
    successors under \texttt{apps/}.
  \item \textbf{Pinned dependencies and environment.} Added
    \texttt{pyproject.toml}, \texttt{requirements.txt}, and an extensive
    \texttt{.env.example} (68 lines) that documents every required key,
    including \texttt{LLM\_BASE\_URL} (defaulted to OpenAI's
    \texttt{api.openai.com/v1}) and the Suno/HealthKit-related settings.
  \item \textbf{Documentation alignment.} The README was rewritten (+357
    lines) to match the new layout, and \texttt{PRODUCT.md} was added to
    state the consumer wellness positioning, anti-references, and design
    principles.
\end{enumerate}

% ============================================================
\section{Music-AI Engine}
% ============================================================

\subsection{Continuous Arousal Modeling}

The mid-term engine mapped physiological signals to a small set of discrete
labels (e.g.\ \texttt{calm}, \texttt{anxious}, \texttt{focus}). This produced
abrupt transitions between recommendations and was unable to express
intermediate states. The final-phase \texttt{processor.py} (rewritten to
\textasciitilde 816 lines) treats arousal as a continuous variable
$a \in [0,1]$ derived from heart-rate variability and EEG-band features
(when available via the MindWave or HealthKit pipeline). Let
$x_t \in \mathbb{R}^d$ be the feature vector at time $t$. The engine produces
a smoothed arousal estimate
\[
  \hat a_t = (1-\lambda)\,\hat a_{t-1} + \lambda\, f_\theta(x_t),
\]
where $f_\theta$ is a feature-to-arousal mapping and $\lambda$ is a smoothing
weight that prevents jitter from short physiological transients.

\subsection{Personalization}

A new module, \texttt{personalization.py} (187 LOC), maintains a lightweight
per-user profile $\pi_u = \langle \text{preferences}, \text{history}, \text{response} \rangle$
and produces a \texttt{MusicStrategy} object for each session:

\[
  \text{Strategy}(u, s) = g\bigl(\hat a_s,\; c_s,\; \pi_u,\; \mathrm{Retrieve}(\hat a_s, c_s)\bigr),
\]

where $s$ denotes the current session, $c_s$ is contextual state (time of day,
self-reported intent), and $\mathrm{Retrieve}(\cdot)$ returns evidence-supported
musical parameters from the knowledge base (Section~\ref{sec:kb}). Strategies
are not opaque prompts; they are structured records carrying tempo (BPM),
mode, instrumentation hints, dynamic envelope, and the IDs of supporting
evidence nodes.

\subsection{Prompt Compilation}

The rewritten \texttt{compiler.py} (361 LOC) transforms a \texttt{MusicStrategy}
into a generation-ready prompt for the downstream music model. We deliberately
separated strategy and prompt so that the strategy layer can be unit-tested
without an external generation service, and so that the same strategy can be
re-targeted at a different generative backend in the future.

% ============================================================
\section{Clinical Music Knowledge Base (GraphRAG)}
\label{sec:kb}
% ============================================================

\subsection{Motivation}

A wellness recommendation must be more than a stylistic guess. To make Soma's
recommendations \emph{auditable}, we built a from-scratch knowledge subsystem
under \texttt{music\_ai\_module/knowledge/} that captures clinical and
musicological evidence as a graph and supports retrieval at strategy time.

\subsection{Data Model}

\texttt{schemas.py} (240 LOC) defines a Pydantic-validated knowledge graph:

\begin{itemize}[leftmargin=*]
  \item \textbf{Entities} --- \emph{Condition} (e.g.\ generalized anxiety),
    \emph{Intervention} (e.g.\ binaural beats at 6\,Hz), \emph{MusicalAttribute}
    (tempo, mode, instrumentation), \emph{Population}, \emph{Outcome}.
  \item \textbf{Relations} --- \texttt{indicates}, \texttt{contraindicates},
    \texttt{modulates}, \texttt{evidenced\_by}.
  \item \textbf{Evidence} --- citations to \texttt{Source} records (paper,
    guideline, dataset) including span / page metadata.
\end{itemize}

\subsection{Pipeline}

\begin{enumerate}[leftmargin=*]
  \item \texttt{fetcher.py} (110 LOC) ingests source documents from configured
    URLs / local paths described in \texttt{data/knowledge/sources.yaml}.
  \item \texttt{extractor.py} (111 LOC) chunks documents and runs LLM-based
    structured extraction into the schemas above.
  \item \texttt{embeddings.py} (100 LOC) computes dense embeddings per chunk
    and per entity surface form.
  \item \texttt{graph\_store.py} (166 LOC) merges new triples into a persistent
    JSON-backed graph (\texttt{data/knowledge/graph.json}, 577 LOC of seed
    data), with deduplication by canonical entity ID.
  \item \texttt{retriever.py} (238 LOC) implements \textbf{hybrid retrieval}.
\end{enumerate}

\subsection{Hybrid Dense--Lexical Retrieval}

For a query $q$ (typically a tuple of $\hat a$, condition tags, and contextual
intent), the retriever scores candidate chunks $d$ using a convex combination:

\[
  \mathrm{score}(d, q) = \alpha \cdot \mathrm{cos}\bigl(\mathbf{e}_d, \mathbf{e}_q\bigr)
                      + (1-\alpha) \cdot \mathrm{BM25}_{\text{multilingual}}(d, q).
\]

The lexical term uses multilingual tokenization so that Chinese clinical
sources and English musicological sources can be retrieved jointly. Top-$k$
chunks are then expanded by a 1-hop graph walk to recover related entities,
and finally passed through \texttt{auditor.py} (152 LOC), which filters out
contraindicated combinations (e.g.\ stimulating tempo for a sleep-induction
session) before the strategy is finalized.

\subsection{Auditing}

The \emph{clinical auditor} is a deterministic post-filter, not a learned
model. It reads the proposed \texttt{MusicStrategy} together with the retrieved
evidence and rejects the strategy if any \texttt{contraindicates} edge is
violated. This gives us a hard safety floor that is easy to test and easy to
explain to a reviewer.

% ============================================================
\section{Web Client Reconstruction}
% ============================================================

The mid-term web client was a single 3{,}834-line HTML file. The final-phase
client (\texttt{apps/web/index.html}, 9{,}866 LOC) is a multi-workspace SPA
organized around a sidebar shell.

\subsection{Architecture}

\begin{itemize}[leftmargin=*]
  \item \textbf{Sidebar App Shell.} A persistent left navigation switches
    between four workspaces: \emph{Session}, \emph{History}, \emph{Diagnostic
    Report}, and \emph{Settings}. Each workspace is independently scrollable
    and uses a unified Tailwind type scale.
  \item \textbf{Single source of truth for vitals.} The \emph{MindWave
    pipeline} is the canonical source for both the displayed vital signs and
    the prompt that is sent to the engine, eliminating an earlier
    inconsistency where the UI display and the model input could drift.
\end{itemize}

\subsection{Session Flow}

The session is split into two screens to reduce cognitive load on entry:

\begin{enumerate}[leftmargin=*]
  \item \textbf{Check-in.} Captures self-report and starts the physiological
    stream; displays a single primary action.
  \item \textbf{Playback.} Plays the generated music with a fixed player and
    independent scroll for parameter context.
\end{enumerate}

We removed the original session-parameter control panel, which had been added
during the prototype phase but interfered with the calm tone the product
intends.

\subsection{Diagnostic Report and Settings}

The diagnostic report defaults to the \emph{Full Report} layout, drops the
unused \emph{Monthly} variant, and uses an LLM-only generation path. A
scrollspy-driven jump navigation was reworked twice to fix overlap with the
report scroll container. The settings page was rebuilt as a sidebar +
tab-panel structure with three sections: \emph{Preferences}, \emph{Privacy},
and \emph{Integrations}.

\subsection{Demo Seeding}

To make the product demonstrable without a live physiological device, a
\texttt{somaSeedDemoHistory()} routine seeds the user's history from five
predefined \texttt{DEMO\_CASES} on first load, and the cached Suno tracks for
those cases were regenerated. Demo report copy was localized to English to
match the audience for the final presentation.

% ============================================================
\section{HealthKit Integration and iOS Foundation}
% ============================================================

A mock HealthKit endpoint, \texttt{POST /api/healthkit/mock/run}, accepts a
JSON payload representing the user's recent HealthKit summary (HRV, resting
HR, sleep stages) and routes it through the same strategy pipeline as
real-time MindWave input. This gave us a complete end-to-end iOS-style flow
without depending on a live device.

In parallel, we drafted
\texttt{docs/superpowers/plans/2026-05-01-ios-migration-foundation.md}
(1{,}782 lines), which lays out the architecture, capability model, data flow,
and migration milestones for moving Soma from a web demo to a SwiftUI iOS
application that consumes HealthKit natively. This document is not part of
the deliverable runtime but defines the path forward for the project.

% ============================================================
\section{Testing}
% ============================================================

The mid-term version had no automated tests. We added 12 modules totaling
approximately 800 lines of test code, summarized in
Table~\ref{tab:tests}.

\begin{table}[h]
\centering
\caption{Test modules added in the final phase.}
\label{tab:tests}
\begin{tabular}{lrl}
\toprule
\textbf{Module} & \textbf{LOC} & \textbf{Subsystem under test}\\
\midrule
\texttt{test\_processor.py}             & 134 & Physiological pipeline \\
\texttt{test\_compiler.py}              & 107 & Prompt compiler \\
\texttt{test\_knowledge\_retriever.py}  & 101 & Hybrid retrieval \\
\texttt{test\_healthkit\_mock\_api.py}  &  88 & FastAPI HealthKit endpoint \\
\texttt{test\_clinical\_auditor.py}     &  75 & Contraindication filter \\
\texttt{test\_pipeline.py}              &  58 & End-to-end orchestration \\
\texttt{test\_knowledge\_schemas.py}    &  54 & Pydantic graph schemas \\
\texttt{test\_knowledge\_seed\_data.py} &  50 & Seed-data integrity \\
\texttt{test\_personalization.py}       &  43 & Per-user strategy \\
\texttt{test\_knowledge\_graph\_merge.py}&  41 & Graph merge / dedup \\
\texttt{test\_knowledge\_extractor.py}  &  16 & Extractor smoke test \\
\texttt{conftest.py} + \texttt{\_\_init\_\_.py} & 32 & Fixtures \\
\midrule
\textbf{Total}                          & \textbf{$\approx$ 800} & \\
\bottomrule
\end{tabular}
\end{table}

The tests run without external network access by relying on local seed data
and a stubbed LLM client, which makes the suite suitable for CI.

% ============================================================
\section{Documentation}
% ============================================================

The final phase materially expanded the documentation set:

\begin{itemize}[leftmargin=*]
  \item \texttt{docs/mvp-prd.md} --- MVP product requirements.
  \item \texttt{docs/stage-1-mvp-session-flow.md} --- session flow specification.
  \item \texttt{docs/session-api-contract.md} --- frontend / backend API contract.
  \item \texttt{docs/demo-script.md} --- demo walkthrough script.
  \item \texttt{docs/superpowers/plans/2026-05-01-ios-migration-foundation.md} ---
        1{,}782-line iOS migration plan.
  \item \texttt{notebooks/pipeline\_demo.ipynb} --- executable engine demo.
\end{itemize}

% ============================================================
\section{Discussion}
% ============================================================

\subsection{What Worked}

The most decisive decision was to pay the cost of the structural refactor
\emph{first}. Once the standard layout, dependency manifest, and test scaffold
were in place, every later subsystem (knowledge base, personalization, API)
could be added without re-litigating import paths, environment loading, or
where things should live. The other strong decision was to separate
\emph{strategy} from \emph{prompt}: by treating the music plan as a typed,
testable record rather than a free-form string, we made the auditor and the
unit tests possible.

\subsection{Limitations}

\begin{itemize}[leftmargin=*]
  \item The clinical knowledge base is seeded but not yet exhaustive; expanding
    coverage requires curated sources and licensed clinical guidelines.
  \item The personalization profile is currently lightweight; longitudinal
    response modeling (does the user actually calm down after sessions of a
    given style?) is future work.
  \item iOS / HealthKit integration is mocked at the API boundary; the native
    SwiftUI implementation is planned but out of scope for this phase.
  \item Continuous arousal estimation depends on signal quality from
    consumer-grade devices and has not been validated against a clinical
    reference.
\end{itemize}

\subsection{Future Work}

Near-term: tighten the auditor's contraindication coverage, expand the
multilingual knowledge sources, and add longitudinal feedback to the
personalization profile. Medium-term: execute the iOS migration plan and
replace the mock HealthKit endpoint with the real SwiftUI bridge.
Longer-term: validate the arousal model against polysomnography or a
clinical-grade HRV reference, and explore a closed-loop session in which the
generation backend is conditioned on real-time physiological response.

% ============================================================
\section{Conclusion}
% ============================================================

Between the mid-term checkpoint and the final submission, Soma evolved from a
working prototype into a structured system: a standard repository layout, a
rewritten music-AI engine with continuous arousal modeling and per-user
strategy, a from-scratch GraphRAG clinical knowledge base with hybrid
multilingual retrieval, a complete SPA reconstruction of the web client, a
mock-HealthKit endpoint, and a test suite covering the core engine and
knowledge subsystems. The combined effect is a product that is no longer
merely able to generate therapeutic music, but able to do so in a way that is
\emph{personalized}, \emph{evidence-grounded}, and \emph{auditable}.

% ============================================================
\appendix
\section{Commit Highlights}
% ============================================================

A non-exhaustive list of the most representative commits in this period:

\begin{longtable}{lp{10cm}}
\toprule
\textbf{Hash} & \textbf{Message}\\
\midrule
\endhead
\texttt{3b0f2f2} & chore: reorganize repo layout (apps, scripts, data, docs, notebooks)\\
\texttt{45e89d7} & feat(music\_ai\_module): continuous arousal model, MusicStrategy, tests\\
\texttt{8079706} & feat(knowledge): add local GraphRAG clinical music knowledge pipeline\\
\texttt{4e9f32e} & Expand knowledge seed data and add integrity tests\\
\texttt{cc61cfa} & knowledge: multilingual lexical tokens and hybrid dense-lexical retrieval\\
\texttt{9ea9287} & feat: personalize music strategy and harden physiology pipeline\\
\texttt{42b4b2b} & feat: mock HealthKit prototype UI and /api/healthkit/mock/run\\
\texttt{b54b6f4} & docs(superpowers): add iOS migration foundation plan\\
\texttt{2c00c1a} & feat(web): Soma MVP session flow with history, feedback, and API contract docs\\
\texttt{9e1f0df} & feat(web): add sidebar app shell with workspace views\\
\texttt{e8433ff} & feat(web): split session into check-in and playback screens\\
\texttt{202292a} & web: unify MindWave pipeline as single source for vitals and prompts\\
\texttt{3dbe611} & Refactor Settings to sidebar layout with tab panels\\
\texttt{94952fc} & polish(web): refine Settings layout, motion curves, and a11y details\\
\texttt{afff047} & feat(web): smooth EQ progress and motion, tighten report flow\\
\bottomrule
\end{longtable}

\end{document}
